\documentclass[a4paper,12pt]{scrbook}
\usepackage{fancyhdr}
\usepackage[a4paper,top=2cm,bottom=2cm]{geometry}
\usepackage{graphicx}
\usepackage[dvipsnames]{pstricks}
\usepackage{pst-text}
\usepackage{pst-grad}
\usepackage{lipsum}

\setlength{\voffset}{0.5in}
\setlength{\oddsidemargin}{28pt}
\setlength{\evensidemargin}{0pt}
\renewcommand{\footrulewidth}{0.5pt}
\lhead[\fancyplain{}{\Large \thepage}]%
      {\fancyplain{}{\rightmark}}
\rhead[\fancyplain{}{\leftmark}]%
      {\fancyplain{}{\Large \thepage}}
\pagestyle{fancy}
\lfoot[\emph{Anto \& Zeli}]{}
\rfoot[]{\emph{Anto \& Zeli}}
\cfoot{}
\topmargin=-0.5in
\textheight=9in
%\sloppy


\makeatletter
\newcommand{\judul}[1]{%
  {\parindent \z@ \centering 
    \interlinepenalty\@M \Large \bfseries #1\par\nobreak \vskip 20\p@ }}
\newcommand{\subjudul}[1]{%
  {\parindent \z@ 
    \interlinepenalty\@M \bfseries #1\par\nobreak \vskip 10\p@ }}
\newcommand{\lagu}[1]{%
  {\parindent \z@ 
    \interlinepenalty\@M \slshape \mdseries \Large \textit{#1}\par\nobreak \vskip 40\p@ }}
\newcommand{\keterangan}[1]{%
  {\parindent \z@  \slshape \large
    \interlinepenalty\@M \textsl{#1}\par\nobreak  \vskip 5\p@}}

\renewenvironment{description}
               {\list{}{\labelwidth\z@ \itemindent-\leftmargin
                        \let\makelabel\descriptionlabel}}
               {\endlist}
\renewcommand*\descriptionlabel[1]{\hspace\labelsep 
                                \normalfont\bfseries #1 }


\makeatother

\newcommand{\BU}[1]{\begin{itemize} \item[U:] #1 \end{itemize}}
\newcommand{\BI}[1]{\begin{itemize} \item[I:] #1 \end{itemize}}
\newcommand{\BP}[1]{\begin{itemize} \item[P:] #1 \end{itemize}}
\newcommand{\BPL}[1]{\begin{itemize} \item[Anto:] #1 \end{itemize}}
\newcommand{\BPW}[1]{\begin{itemize} \item[Zeli:] #1 \end{itemize}}
\newcommand{\BPLW}[1]{\begin{itemize} \item[A+Z:] #1 \end{itemize}}
\newcommand{\BOT}[1]{\begin{itemize} \item[OT:] #1 \end{itemize}}
\newcommand{\BS}[1]{\begin{itemize} \item[Saksi:] #1 \end{itemize}}
\newcommand{\BW}[1]{\begin{itemize} \item[Wakil:] #1 \end{itemize}}
\newcommand{\bpcamantri}{Yohanes Subaryadi }
\newcommand{\ibucamantri}{Veronika Suwardiyanti }
\newcommand{\bpcamantra}{Yohanes Sudarmadi }
\newcommand{\ibucamantra}{Anastasia Djuwarni }
\newcommand{\saksisatu}{Bapak Drs. B. Kusmanto  }
\newcommand{\saksidua}{Bapak Drs. Neo Suradi  }

\newcommand{\romo}{Robertus Triwidodo, Pr }
\newcommand{\camantri}{R Zeli Puspitasari}
\newcommand{\camantra}{FX Arie Wibowo Sudaryanto}

\newcommand{\namagereja}{Gereja Bunda Maria Maguwo }

%lagu-lagu
\newcommand{\lagupembukaan}{Aku Bersyukur (\textsl{Asri Savsavubun/koor}) \\ \\

Di hari ini aku bersyukur, atas rahmatMu Tuhan. \\
Di hari ini aku bahagia, atas rahmatMu Tuhan. \\ 
Syukur pujian kusampaikan kepada Tuhan \\
Tuhan selalu membimbingku syukur pada Tuhan
}
\newcommand{\lagukyrie}{MB 186 }
\newcommand{\lagugloria}{~}
\newcommand{\laguantarbacaan}{Hatiku (koor) \\ \\
Reff:\\
Hatiku kusembahkan padaMu\\
Kasih abadi Kaulimpahkan kepadaku\\ \\
HatiMu slalu murah dan penyayang \\
Hatiku slalu cenderung berdosa\\
Tapi cintaMu mengiring langkahku\\
Yesus penebusku \\ \\
KuasaMu slalu memeluk diriku\\
Kuasaku slalu menentang hukumMu\\
Tapi cintaMu mengiring langkahku\\
Yesus penebusku \\ \\
JalanMu slalu mengarah ke Bapa\\
Jalanku sering diliput gulita\\
Tapi cintaMu mengiring langkahku\\
Yesus penebusku
}
\newcommand{\lagubaitpengantarinjil}{~}
\newcommand{\lagupasangcincin}{Cincin Kami (\textsl{Esin Savsavubun}) \\ \\
Cincin tanda kami cinta kasih sejati, tlah melingkar di jemari.\\
Cincin tanda kami cinta kasih sejati, tlah mengikat dua hati. \\ 

Di hadapan altar yang suci ini, kita saling serah diri. \\
Di hadapan altar yang suci ini, cinta kita kan abadi. \\

Badai taufan kan menghadang, namun cinta kita tak kan goyang.\\
Tuhan kan memberi jalan serta cahayanya mengarungi samodera cinta.\\

Tuhan bimbinglah kami berdua berlayar dalam bahtera.\\
Tuhan bimbinglah kami berdua hidup seturut sabdaMu.\\

Badai taufan kan menghadang, namun cinta kita tak kan goyang.\\
Tuhan beri kami kasih serta cinta, spanjang hdup kami.\\
}
\newcommand{\lagusungkem}{Di Doa Ibuku Namaku Disebut (\textsl{Asri
Savsavubun})\\

Di waktu ku kecil, gembira dan senang\\
Tiada duka kukenal, tak kunjung mengerang\\
Disore hari nan sepi....ibuku terbelut\\
Sujud berdoa ku dengar namaku disebut\\

Di doa ibuku, namaku disebut\\
Di doa ibuku kudengar, ada namaku disebut\\

Sering ini kukenang, di masa yang berat\\
Di kala hidup mendesak dan nyaris kutersesat\\
Melintas gambar ibuku, sewaktu bertelut\\
Kembali sayup kudengar, ....namaku disebut\\

Di sore hari nan sepi... ibuku bertelut\\
Sujud berdoa ku dengar namaku disebut\\
Di doa ibuku, namaku disebut\\
Di doa ibuku dengar ada namaku disebut....\\
Ada namaku di sebut }

\newcommand{\lagusungkemmaria}{Ave Maria (\textsl{Esin Savsavubun})\\

Ave Maria Gratia plena\\
Maria Gratia plena\\
Maria Gratia plena\\
Ave, ave dominus\\
Dominus tecum\\

Benedicta tu in mulieribus\\
Et benedictus\\
Et benedictus fructus ventris\\
Ventris tui Jesus\\
Ave Maria\\
}

\newcommand{\lagupersembahan}{Betapa Hatiku (koor) \\

Betapa hatiku berterimakasih Yesus \\
Kau mengasihiku, Kau memilikiku, \\

Hanya ini Tuhan persembahanku \\
Segenap hidupku jiwa dan ragaku\\
Sbab tak kumiliki harta kekayaan\\
Yang cukup berarti tuk kupersembahkan \\

Hanya ini Tuhan permohonanku\\
Terimalah Tuhan persembahanku\\
Pakailah hidupku sebagai alatMu, seumur hidupku
}
\newcommand{\lagusanctus}{~ }
\newcommand{\lagubapakami}{MB 144 }
\newcommand{\laguagnusdei}{MB 278 }
\newcommand{\lagukomuni}{Kumau Cinta Yesus (\textsl{Asri Savsavubun})\\

Kumau cinta Yesus selamanya\\
Kumau cinta Yesus selamanya\\
Meskipun badai silih berganti dalam hidupku\\
Ku tetap cinta Yesus selamanya\\

Ya Bapa Bapa, ini aku anakMu
Layakkanlah seluruh hidupku
Ya Bapa Bapa, ini aku anakMu
Pakailah sesuai dengan rencanaMu
}
\newcommand{\lagupenutup}{O Sanctisima (koor)\\

    O Sanctissima, O Piissima\\
    Dulcis Virgo Maria\\
    Mater amata, Intemerata\\
    Ora, Ora Pro Nobis\\

Tu solatium et refugium\\
Virgo Mater Maria\\
Quidquid optamus per te speramus\\
Ora, ora pro nobis
}


\usepackage[bahasa]{babel}
\selectlanguage{bahasa}
\hyphenation{sa-u-da-ra-ku}
\hyphenation{ke-ri-ngat}
\hyphenation{je-ri-tan}
\hyphenation{hu-bung-an}
\hyphenation{me-nya-dari}
\hyphenation{Eng-kau}
\hyphenation{ke-sa-lah-an}
\hyphenation{ba-gai-ma-na}
\hyphenation{Tu-han}
\hyphenation{di-per-ca-ya-kan}
\hyphenation{men-ja-uh-kan}
\hyphenation{bu-kan-lah}
\hyphenation{per-sa-tu-kan-lah}
\hyphenation{ma-khluk}
\hyphenation{Sem-buh-kan-lah}
\hyphenation{ja-lan}
\hyphenation{mem-bu-tuh-kan}
\hyphenation{be-ri-kan-lah}
\hyphenation{me-ra-sa-kan}
\hyphenation{te-man-ilah}
\hyphenation{mem-bi-ngung-kan}
\hyphenation{di-ka-gum-i}
\hyphenation{ta-ngis-an-Mu}
\hyphenation{mi-lik-ilah}


\title{EKARISTI \\SAKRAMEN PERKAWINAN}
\author{
\includegraphics[scale=0.25]{anto-zeli.eps}\\ \\
\camantra ~dan \camantri \\
\\ \\
oleh \\Romo \romo \\
\\
\namagereja} 
\date{13 Juni 2009}

\begin{document}
\thispagestyle{empty}
\newsavebox\IBox
\sbox\IBox{\includegraphics[scale=0.25]{anto-zeli.eps}}
% set up the picture environment
\psset{unit=1in}
\begin{pspicture}(5in,8in)
% set up the fonts we use
\DeclareFixedFont{\PT}{T1}{ppl}{b}{it}{0.5in}
\DeclareFixedFont{\PTsmall}{T1}{ppl}{b}{it}{0.4in}
\DeclareFixedFont{\PTsmallest}{T1}{ppl}{b}{it}{0.3in}
\DeclareFixedFont{\PTtext}{T1}{ppl}{b}{it}{11pt}
\DeclareFixedFont{\Logo}{T1}{pbk}{m}{n}{0.3in}
% place the front cover picture
\rput[cb](3,4){\usebox\IBox}
% put the text on the front cover
\rput[cb](3,7){\PT {EKARISTI}}
\rput[cb](3,6.25){\PT {SAKRAMEN PERKAWINAN}}
\rput[cb](3,2.2){\PTsmallest {\camantra}}
\rput[cb](3,1.8){\PTsmallest {dan}}
\rput[cb](3,1.4){\PTsmallest {\camantri }}
\rput[cb](3,0.4){\PTsmallest {oleh}} 
\rput[cb](3,0){\PTsmallest {Romo \romo}}

%\rput[cb](3,-1){\PTsmallest {\namagereja}} 

\end{pspicture}

\begin{center}\Large \namagereja

13 Juni 2009
\end{center}

\sffamily
%\maketitle
\Large
\newpage
\thispagestyle{empty}
{~}
\newpage
\thispagestyle{empty}

\judul{UPACARA PENYAMBUTAN}
\keterangan{Kedua mempelai bersiap di pintu utama gereja didampingi oleh orang
tua dan keluarga, serta salah seorang wakil keluarga untuk menyerahkan mempelai
kepada gereja.}

\BP{Selamat datang Bapak, Ibu, dan Saudara-saudara sekalian di \namagereja ini.
Pada hari ini kita akan memberikan dukungan dan dorongan doa atas Penerimaan
Sakramen Perkawinan antara \camantra ~putra Bapak \bpcamantra dan Ibu
\ibucamantra dengan \camantri ~puteri Bapak \bpcamantri dan Ibu \ibucamantri yang
akan dipimpin oleh Romo \romo.

Allah menciptakan manusia menurut citra-Nya sebagai pria dan wanita. Mereka
diberi kemampuan untuk saling mencintai dan Allah memberkati serta berfirman:
'Beranak cuculah dan bertambah banyak dan penuhilah bumi dan taklukkanlah'.

Dengan demikian perkawinan oleh Allah sebagai suatu lembaga keselamatan sebab
dengan membangun keluarga manusia dapat menghayati darma baktinya kepada Allah
demi kelangsungan dunia, yaitu melaksanakan cinta kasih, daya cipta dan kuasa
mengatur yang berasal dari Allah.

Acara ini akan diawali dengan penyambutan kedua mempelai oleh Romo, marilah kita
berdiri.}

\keterangan{Imam dan putera altar menyambut mempelai di depan pintu}

\subjudul{Penerimaan Mempelai}
\BI{Saudara sekalian yang terkasih dalam Kristus, selamat datang di rumah Tuhan
ini. Semoga damai dan kasih Allah selalu besertamu.}
\BU{Dan sertamu juga.}

\subjudul{Penyerahan Mempelai}
\BW{Romo yang terhormat, perkenankanlah saya, selaku wakil keluarga mempelai,
menyerahkan putra putri kami \camantra dan \camantri ke hadapan Romo \romo agar
Romo berkenan mengukuhkan pernikahan mereka menurut tata cara Gereja Katolik
yang kudus.}

\BI{Selaku pejabat gereja katolik, saya mengucapkan terima kasih, dan dengan
tulus ikhlas menerima kedua mempelai ini. Saya akan mohonkan berkat Tuhan agar
ikatan perkawinan mereka sungguh diteguhkan seumur hidup.
Dan semoga air suci ini menyucikan saudara-saudari agar pantas menerima berkat
Tuhan dalam upacara yang agung ini.

\keterangan{(Imam memerciki mempelai dan rombongan dengan air suci)}
Dan sekarang saya mengundang kedua mempelai bersama seluruh keluarga untuk
menuju altar Tuhan.}


\subjudul{Perarakan} 
\keterangan{Umat berdiri, Romo, Putra Altar, kedua mempelai, para saksi serta
orang tua kedua mempelai berarak menuju ke depan altar.}

\lagu{Lagu pembukaan - \lagupembukaan}

\judul{RITUS PEMBUKA}

\subjudul{Tanda Salib dan Salam}
\BI{Demi Nama Bapa dan Putera dan Roh Kudus}
\BU{Amin}
\BI{Kasih karunia, rahmat dan damai sejahtera dari Allah Bapa dan dari
Putera-Nya Yesus Kristus, besertamu}
\BU{Dan sertamu juga}

\subjudul{Kata Pengantar dari Romo}

\subjudul{Pernyataan Tobat}
\BI{Saudara saudari terkasih, marilah kita menyesali segala dosa-dosa kita dan
mohon ampun agar kita layak merayakan perjamuan kudus ini.}
\BI{Saya mengaku \dots}
\BU{Kepada Allah yang Maha Kuasa dan kepada saudara sekalian bahwa saya telah
berdosa dengan pikiran dan perkataan, dengan perbuatan dan kelalaian. Saya
berdosa, saya berdosa, saya sungguh berdosa. Oleh sebab itu saya mohon kepada
Santa Perawan Maria, kepada para malaikat dan orang kudus, dan kepada saudara
sekalian supaya mendoakan saya kepada Allah Tuhan kita.}
\BI{Semoga Allah yang Maha Kuasa mengasihi kita, mengampuni dosa kita, dan
menghantar kita ke hidup yang kekal.}
\BU{Amin}

\lagu{Tuhan Kasihanilah Kami \lagukyrie }

\subjudul{Doa Pembukaan}
\BI{Marilah kita berdoa

Allah pencipta dan penebus kami, Engkau menghendaki agar pria dan wanita
membangun keluarga yang bahagia. Kedua hamba-Mu ini: \camantra ~dan \camantri
~sudah siap untuk memasuki bahtera perkawinan.

Berkatilah cinta kasih mereka, supaya tahan uji dalam untung dan malang.
Anugerahkanlah mereka berdua kesetiaan, kebahagiaan, serta keturunan yang dapat
dibanggakan. Demi Yesus Kristus, Putra-Mu, Tuhan dan pengantara kami, yang hidup
dan berkuasa kini dan sepanjang masa.}

\BU{Amin.}

\judul{LITURGI SABDA}
\subjudul{Bacaan pertama Pengkhotbah 3: 1 - 11}
\BP{ 
\textit{Untuk segala sesuatu ada waktunya}

Untuk segala sesuatu ada masanya, untuk apapun di bawah langit ada waktunya.
Ada waktu untuk lahir, ada waktu untuk meninggal, ada waktu untuk menanam, ada waktu untuk mencabut yang ditanam;
ada waktu untuk membunuh, ada waktu untuk menyembuhkan; ada waktu untuk merombak, ada waktu untuk membangun;
ada waktu untuk menangis, ada waktu untuk tertawa; ada waktu untuk meratap; ada waktu untuk menari;
ada waktu untuk membuang batu, ada waktu untuk mengumpulkan batu; ada waktu untuk memeluk, ada waktu untuk menahan diri dari memeluk;
ada waktu untuk mencari, ada waktu untuk membiarkan rugi; ada waktu untuk menyimpan, ada waktu untuk membuang;
ada waktu untuk merobek, ada waktu untuk menjahit; ada waktu untuk berdiam diri, ada waktu untuk berbicara;
ada waktu untuk mengasihi, ada waktu untuk membenci; ada waktu untuk perang, ada waktu untuk damai.

Apakah untung pekerja dari yang dikerjakannya dengan berjerih payah?

Aku telah melihat pekerjaan yang diberikan Allah kepada anak-anak manusia untuk melelahkan dirinya.
Ia membuat segala sesuatu indah pada waktunya, bahkan Ia memberikan kekekalan dalam hati mereka. Tetapi manusia tidak dapat menyelami pekerjaan yang dilakukan Allah dari awal sampai akhir.

Demikianlah sabda Tuhan}

\BU{Syukur kepada Allah}

\lagu{Lagu Pengantar Bacaan - \laguantarbacaan}

\subjudul{Bacaan Injil}
\BI{Tuhan sertamu}
\BU{Dan sertamu juga}
\BI{Inilah Injil Suci Yesus Kristus menurut Santo Yohanes (Yoh 2: 1-10)}
\BU{Terpujilah Kristus}

\BI{
Pada hari ketiga ada perkawinan di Kana yang di Galilea, dan ibu Yesus ada di situ;
Yesus dan murid-murid-Nya diundang juga ke perkawinan itu.
Ketika mereka kekurangan anggur, ibu Yesus berkata kepada-Nya: "Mereka kehabisan anggur."
Kata Yesus kepadanya: "Mau apakah engkau dari-Ku, ibu? Saat-Ku belum tiba."
Tetapi ibu Yesus berkata kepada pelayan-pelayan: "Apa yang dikatakan kepadamu, buatlah itu!"

Di situ ada enam tempayan yang disediakan untuk pembasuhan menurut adat orang Yahudi, masing-masing isinya dua tiga buyung.
Yesus berkata kepada pelayan-pelayan itu: "Isilah tempayan-tempayan itu penuh dengan air." Dan mereka pun mengisinya sampai penuh.
Lalu kata Yesus kepada mereka: "Sekarang cedoklah dan bawalah kepada pemimpin pesta." Lalu mereka pun membawanya.

Setelah pemimpin pesta itu mengecap air, yang telah menjadi anggur itu -- dan ia tidak tahu dari mana datangnya, tetapi pelayan-pelayan, yang mencedok air itu, mengetahuinya -- ia memanggil mempelai laki-laki,
dan berkata kepadanya: "Setiap orang menghidangkan anggur yang baik dahulu dan sesudah orang puas minum, barulah yang kurang baik; akan tetapi engkau menyimpan anggur yang baik sampai sekarang."

Demikianlah Injil Tuhan}

\BU{Terpujilah Kristus}

\subjudul{Homili}

\judul{LITURGI PERKAWINAN}

\subjudul{Pernyataan Saksi}
\keterangan{Mempelai berdiri didampingi oleh kedua saksi dan orang tua}

\BI{Mempelai berdua yang berbahagia, saudara telah hadir di gereja ini untuk
melangsungkan perkawinan, di hadapan pejabat Gereja dan disaksikan oleh umat
beriman. Kiranya Tuhan memberkati dan meneguhkan saudara. Namun sebelum
perkawinan ini diresmikan perkenankan saya meminta pernyataan dari para saksi.
Para saksi yang terhormat, adakah sesuatu yang menghalangi pernikahan ini
menurut adat dan ajaran Gereja Katolik?}

\BS{Sepanjang pengetahuan kami, tidak ada halangan apapun untuk meresmikan
pernikahan ini menurut tata cara gereja katolik. Kami mendukung sepenuhnya
permohonan kedua mempelai, dan kami bersedia menjadi saksinya.}

\subjudul{Pernyataan Mempelai}
\BI{Mempelai berdua yang berbahagia, setelah saya mengadakan penyelidikan
seperlunya, saya selaku pejabat Gereja meluluskan permintaan Anda berdua. Akan
tetapi sebelum perkawinan Anda berdua ini diresmikan, saya persilahkan kedua
mempelai menyatakan kesungguhan hati Anda berdua di hadapan Allah dan semua yang
hadir.}

\BPL{Saya, \camantra, menyatakan bahwa saya sungguh ingin meresmikan perkawinan
ini dengan ikhlas hati. Saya bersedia mencintai dan menghormati istri saya
sepanjang hidup saya. Saya bersedia menjadi bapak yang baik, bagi anak-anak yang
dipercayakan Tuhan kepada saya serta mendidik mereka dalam terang iman
Kristiani.}

\BPW{Saya, \camantri, menyatakan bahwa saya sungguh ingin meresmikan perkawinan
ini dengan ikhlas hati. Saya bersedia mencintai dan menghormati suami saya
sepanjang hidup saya. Saya bersedia menjadi ibu yang baik, bagi anak-anak yang
dipercayakan Tuhan kepada saya serta mendidik mereka dalam terang iman
Kristiani.}

\subjudul{Perjanjian Perkawinan}
\BI{Maka tibalah saatnya untuk meresmikan perkawinan Anda berdua. Saya
mempersilahkan Anda masing-masing mengucapkan perjanjian perkawinan satu-persatu
sambil meletakkan tangan di atas Kitab Suci ini. Dengan harapan janji saudara
berdua menjadi ikatan seumur hidup dalam kasih Tuhan.}

\subjudul{Janji Mempelai Pria}
\BPL{Di hadapan Romo dan para saksi, saya \camantra ~menyatakan dengan tulus
ikhlas bahwa \camantri ~yang hadir di sini, mulai sekarang ini menjadi istri
saya. Saya berjanji setia kepadanya dalam untung dan malang dan saya mau
mencintai dan menghormatinya seumur hidup. Demikianlah janji saya, demi Allah
dan Injil suci ini.}

\subjudul{Janji Mempelai Wanita}
\BPW{Di hadapan romo dan para saksi, saya \camantri ~menyatakan dengan tulus
ikhlas bahwa \camantra ~yang hadir di sini, mulai sekarang ini menjadi suami
saya. Saya berjanji setia kepadanya dalam untung dan malang dan saya mau
mencintai dan menghormatinya seumur hidup. Demikianlah janji saya, demi Allah
dan Injil suci ini.}

\subjudul{Peneguhan Perkawinan}
\BI{Atas nama Gereja Allah dan di hadapan para saksi dan hadirin sekalian, saya
menegaskan bahwa perkawinan yang telah diresmikan ini adalah perkawinan Katolik
yang sah. Semoga sakramen ini menjadi bagi Anda berdua sumber kekuatan dan
kebahagiaan. Yang dipersatukan Allah, janganlah diceraikan manusia.}
\BU{Amin}

\subjudul{Doa untuk Mempelai}
\BI{Saudara saudari yang terkasih, marilah kita berdoa dengan rendah hati,
supaya Tuhan rela mencurahkan berkat dan rahmat-Nya atas kedua mempelai ini yang
telah menikah dalam Kristus. Mereka sudah dipersatukan dalam perjanjian suci.
Semoga mereka bersatu padu dalam cinta kasih}

\BI{Allah, Engkau menciptakan segala sesuatu dengan kekuatan kuasa-Mu. Engkau
menciptakan manusia menurut citra-Mu. Engkau menciptakan pria dan wanita supaya
mereka dipadukan menjadi satu. Engkau mengajarkan bahwa perkawinan yang telah
Kau teguhkan tak boleh diceraikan.
Engkau juga menguduskan ikatan perkawinan dan mengangkat perjanjian perkawinan
menjadi lambang persatuan Kristus dengan gereja-Nya.

Pandanglah dengan rela mempelai wanita ini agar rahmat cinta dan damai sejahtera
tinggal dalam hatinya. Semoga ia menjadi istri yang setia dan ibu yang baik
seperti wanita-wanita yang kudus yang dipuji dalam Kitab Suci.

Kami berdoa pula untuk mempelai pria ini semoga ia selalu berusaha untuk
mencintai dengan tulus ikhlas dan menunaikan tanggung jawabnya baik terhadap
istri dan anak-anak maupun terhadap masyarakat.

Dan kini kami mohon kepada Mu ya Tuhan,
Semoga kedua mempelai ini tetap berpegang pada iman dan perintah-perintah-Mu.
Semoga mereka bersatu sebagai suami-istri, terpandang karena peri hidup yang
baik dan berjasa untuk sesama dalam lingkungan mereka. Kuatkanlah mereka dengan
semangat Injil sehingga menjadi saksi Kristus bagi semua orang. Kuatkanlah
mereka dengan Roh Kudus-Mu agar mampu menghadapi setiap godaan dan tantangan
kehidupan. Dan akhirnya mereka mampu pula menjadi keluarga yang berbahagia.
Semoga mereka subur dan berketurunan menjadi orangtua yang patut dicontoh dan
berbahagia melihat anak cucunya. Semoga mereka mencapai usia lanjut dan akhirnya
memasuki kehidupan bahagia dalam kerajaan surga.

Demi Kristus, pengantara kami.}

\BU{Amin}

\subjudul{Pemberkatan Cincin}
\keterangan{Diiringi lagu \lagupasangcincin}
\BI{Ya Allah sumber kesetiaan, berkatilah kedua cincin ini supaya menjadi
lambang cinta dan kesetiaan tanpa akhir. Semoga kedua cincin ini mengingatkan
kedua mempelai akan janji setia yang mereka ucapkan pada hari ini. Demi Kristus
Tuhan dan pengantara kami.}

\keterangan{Cincin diserahkan kepada mempelai pria}

\BI{Kenakanlah cincin ini pada jari istrimu sebagai lambang cinta dan
kesetiaanmu padanya.}
\BPL{\camantri, terimalah cincin ini sebagai tanda cinta kasihku kepadamu.}

\keterangan{Cincin diserahkan kepada mempelai wanita}

\BI{Kenakanlah cincin ini pada jari suamimu sebagai lambang cinta dan
kesetiaanmu padanya.}
\BPW{\camantra, terimalah cincin ini sebagai tanda cinta kasihku kepadamu.}

\lagu{Mohon Doa Restu Orantua: lagu \lagusungkem}

\subjudul{Pemberkatan Salib, Kitab Suci, dan Rosario}

\BI{Ya Tuhan berkatilah Kitab Suci ini agar dapat dipakai oleh keluarga baru ini
sebagai sumber pengharapan, penghiburan, dan cinta kasih sejati dalam hidup
mereka.
Berkatilah pula Salib dan Rosario ini, agar selalu menjadi tanda kehadiran-Mu
dan Bunda Maria. Dengan pengantaraan Kristus, Tuhan kami.} 
\BU{Amin.}

\keterangan{Romo menyerahkan Salib, Kitab Suci, dan Rosario kepada Putra Altar
yang kemudian menyerahkan kepada wakil orang tua. Orang tua menyerahkan kepada
kedua mempelai.}

\BOT{Sebagai orangtua kalian berdua, saya menyerahkan kitab suci, salib, dan
rosario ini, sebagai lambang bahwa kami mewariskan kepadamu berdua harta yang
tak ternilai ini untuk bekal kalian dalam perjalanan hidupmu, yakni iman akan
kasih Tuhan Yesus yang telah mengorbankan diriNya.

Semoga iman akan Salib dan Firman Tuhan memberi harapan dan penghiburan hati
kalian di kala duka dan membangkitkan semangat cinta sejati di saat gersang.
Semoga berkat Rosario yang kudus ini, kalian berdua akan selalu dekat dengan
Bunda Maria, Ibu kita, dan Ibu setiap keluarga Kristus.}

\BPLW{Amin.}

\subjudul{Penandatanganan Dokumen Pernikahan}
\BI{Saudara sekalian, kita telah menyaksikan peresmian pernikahan kedua mempelai
ini. Maka sekarang saya persilahkan kedua mempelai dan para saksi untuk
menandatangani dokumen pernikahan.}

\keterangan{Romo, para saksi, dan kedua mempelai menandatangani berkas
pernikahan.
Sesudah itu mempelai kembali ke tempat duduk.}

\subjudul{Doa Umat}
\BI{Saudara saudari yang terkasih, kedua mempelai ini telah mengikat perjanjian
nikah di hadapan Allah dan di hadapan para hadirin. Marilah kita mendoakan
mereka agar mereka tetap bahagia dalam lindungan Tuhan.}

\BP{Semoga mereka menghayati hidup perkawinan dalam cinta kasih dan damai,
sehingga rahmat dan kebaikan Kristus bersinar dalam rumah tangga mereka. Kami
mohon..}
\BU{Kabulkanlah doa kami}

\BP{Semoga mereka sungguh mendasarkan hidup perkawinan atas dasar cinta kasih
sejati, sebagaimana cinta Yesus kepada gereja-Nya. Kami mohon..}
\BU{Kabulkanlah doa kami}

\BP{Semoga cinta kasih mereka diberkati oleh Tuhan dengan karunia yang
berlimpah, dengan kemurahan rezeki sehari-hari, dengan kesehatan, dengan
kerukunan dan damai; sehingga anak-anak yang dianugerahkan kepada mereka
sungguh-sungguh menggembirakan hati orang tuanya.
Kami mohon..}
\BU{Kabulkanlah doa kami}

\BP{Semoga orang tua mereka merasa terhibur karena menyaksikan kebahagiaan
anaknya, sehingga pengantin baru ini sungguh diterima dalam lingkungan kaum
kerabatnya.
Kami mohon..}
\BU{Kabulkanlah doa kami}

\BI{Ya Tuhan, peliharalah saudara-saudari kami ini dalam cinta-Mu, bantulah dan
tuntunlah mereka seumur hidup. Dan kabulkanlah doa-doa kami. Demi Kristus Tuhan
dan Pengantara kami.}


\judul{LITURGI EKARISTI}

\subjudul{Perarakan Persembahan}
\lagu{Lagu Persembahan - \lagupersembahan}

\subjudul{Doa Persembahan}
\BI{Allah Bapa kami di surga, terimalah persembahan yang kami hunjukkan untuk
menyucikan perkawinan ini. Semoga mereka menikmati berkat-Mu dalam membangun
keluarga. Demi Kristus pengantara kami}
\BU{Amin}

\subjudul{Prefasi}
\BI{Tuhan sertamu}
\BU{Dan sertamu juga}
\BI{Marilah mengarahkan hati kepada Tuhan}
\BU{Sudah kami arahkan}
\BI{Marilah bersyukur kepada Tuhan Allah kita}
\BU{Sudah layak dan sepantasnya}
\BI{Sungguh layak dan sepantasnya ya Bapa yang Kudus Allah yang Kekal dan Kuasa,
bahwa di mana pun juga kami senantiasa bersyukur kepada-Mu sebab Engkau berkenan
menciptakan manusia dan menganugerahinya martabat yang luhur. Persatuan mesra
antara suami istri mencerminkan cinta kasih-Mu sendiri. Engkau menciptakan
manusia karena cinta. Engkau memanggil dia untuk mengamalkan cinta. Engkau
mengundang dia untuk menikmati cinta-Mu yang kekal. Dengan demikikan misteri
perkawinan suci menandakan cinta Ilahi dan menguduskan cinta insani. Maka
bersama para malaikat dan seluruh rombongan para kudus kami bermadah memuliakan
Dikau dengan tak henti-hentinya berseru:}

\lagu{Kudus}

\subjudul{Doa Syukur Agung II}

\BI{Sunguh kuduslah Engkau ya Bapa. Segala ciptaan patut memuji Engkau. Sebab
dengan pengantaraan Putra-Mu, Tuhan kami Yesus Kristus, dan dengan daya kekuatan
Roh Kudus Engkau menghidupkan dan menguduskan segala sesuatu. Tak henti-hentinya
Engkau menghimpun umat-Mu, sehingga dari terbitnya matahari sampai terbenamnya
di seluruh bumi dipersembahkan kurban yang murni untuk memuliakan nama-Mu. Maka
kami mohon ya Bapa sudilah menguduskan persembahan ini dengan Roh-Mu agar bagi
kami menjadi tubuh ($\dagger$) dan darah Putra-Mu terkasih Tuhan kami, Yesus
Kristus yang menghendaki kami merayakan misteri ini.

Sebab pada malam ia dikhianati, Yesus mengambil roti, Ia mengucap syukur dan
memuji Dikau, memecah-mecahkan roti itu dan memberikannya kepada
murid-murid-Nya seraya berkata:

TERIMALAH DAN MAKANLAH! INILAH TUBUHKU, YANG DISERAHKAN BAGIMU!

Demikian pula sesudah perjamuan, Yesus mengambil piala. Sekali lagi Ia mengucap
syukur dan memuji Dikau, lalu memberikan piala itu kepada murid-murid-Nya seraya
berkata:

TERIMALAH DAN MINUMLAH INILAH PIALA DARAHKU, DARAH PERJANJIAN BARU DAN KEKAL,
YANG DITUMPAHKAN BAGIMU DAN BAGI SEMUA ORANG DEMI PENGAMPUNAN DOSA. LAKUKANLAH
INI UNTUK MENGENANGKAN DAKU!}

\lagu{Anamnese}

\BI{Bapa, kami mengenangkan sengsara Putra-Mu, yang menyelamatkan,
kebangkitan-Nya yang mengagumkan, dan kenaikkan-Nya ke surga. Sambil
mengharapkan kedatangan-Nya kembali, dengan penuh syukur kami mempersembahkan
kepada-Mu kurban yang hidup dan kudus ini. Kami mohon, pandanglah persembahan
Gereja-Mu ini dan indahkanlah kurban yang telah mendamaikan kami dengan Dikau
ini.

Semoga kami disempurnakan oleh-Nya menjadi suatu persembahan abadi bagi-Mu, agar
kami pantas mewarisi kebahagiaan surgawi bersama para pilihan-Mu, terutama
bersama Santa Perawan Maria, Bunda Allah, para rasul-Mu yang kudus, dan para
martir-Mu yang jaya, dan bersama semua orang kudus yang selalu mendampingi dan
menolong kami.

Ya Bapa, semoga berkat korban yang mendamaikan ini, damai sejahtera dan
keselamatan semakin dirasakan oleh dunia.

Kuatkanlah Iman dan cinta kasih Gereja-Mu yang kini masih berziarah di bumi ini
bersama hamba-Mu, Paus Benedictus XVI dan Uskup kami Ignatius Soeharyo, serta
semua uskup, para imam, diakon serta semua pelayan umat, dan seluruh umat
kesayangan-Mu. Dengarkanlah doa-doa umat-Mu yang berhimpun di sini. Demi
kerahiman dan kasih setia-Mu, ya Bapa, persatukanlah semua anak-Mu di mana pun
mereka berada.

Sudilah pula menganugerahkan kebahagiaan abadi kepada semua yang telah berpulang
ke hadirat-Mu, saudara-saudara kami seiman, dan semua orang lain yang hidupnya
berkenan pada-Mu. Pada waktu itu, Engkau menghapus setiap tetes air mata kami,
karena dengan memandang Engkau, ya Bapa, kami akan serupa dengan Dikau sepanjang
masa dan tak henti-hentinya memuji Dikau.

Kami berharap, agar bersama mereka kami pun menikmati kemuliaan-Mu
selama-lamanya, dengan pengantsraan Kristus, Tuhan kami. Sebab melalui Dialah
Engkau melimpahkan segala yang baik kepada dunia.

Dengan pengantaraan Kristus, bersama Dia dan dalam Dia, bagi-Mu Allah Bapa Yang
Maha Kuasa, dalam persekutuan dengan Roh Kudus, segala hormat dan kemuliaan,
sepanjang segala mas. Amin.}

\subjudul{Bapa Kami}
\BI{Atas petunjuk Penyelamat kita dan menurut ajaran Ilahi, maka beranilah
kita berdoa}

\BU{Bapa Kami yang ada di surga \dots}

\subjudul{Doa Damai}
\BI{Tuhan Yesus pernah bersabda kepada para rasul: “Damai Kuwartakan kepadamu”.
Maka marilah kita berdoa kepada-Nya:}

\BU{Tuhan Yesus Kristus, janganlah memperhitungkan dosa kami, tetapi
perhatikanlah iman Gereja-Mu, dan restuilah kami, khususnya kedua mempelai ini,
supaya hidup bersatu dengan rukun sesuai dengan kehendak-Mu. Sebab Engkaulah
pengantara kami kini dan sepanjang masa. Amin.}
\BI{Semoga damai Tuhan kita Yesus Kristus selalu besertamu.}
\BU{Dan sertamu juga.}


\subjudul{Anak Domba Allah}
\BI{Inilah Anak Domba Allah yang menghapus dosa dunia, berbahagialah saudara
yang diundang ke perjamuan-Nya}
\BU{Ya Tuhan, saya tidak pantas Tuhan datang pada saya, tetapi bersabdalah saja,
maka saya akan sembuh.}

\subjudul{Komuni}
\lagu{Lagu Komuni - \lagukomuni}
\keterangan{Umat yang telah dibabtis secara Katolik maju untuk menerima komuni}


\subjudul{Doa Penutup}
\BI{Allah yang Maha Kuasa, kami mohon semoga kedua mempelai yang berbahagia ini
semakin diresapi oleh rahmat sakramen perkawinan yang telah kami rayakan dan
kuatkanlah kami semua dengan daya ekaristi suci supaya kami lebih sanggup hidup
sebagai murid Kristus yang sejati. Dialah pengantara kami sepanjang segala
masa.}
\BU{Amin.}

\subjudul{Ucapan Terima Kasih dari Keluarga}

\subjudul{Berkat Penutup}
\BI{Tuhan sertamu}
\BU{Dan sertamu juga}
\BI{Saudara-saudara, marilah kita akhiri perayaan ini dengan memohon berkat
Tuhan, khususnya bagi kedua mempelai}
\BI{Semoga Allah Bapa yang mahakuasa menganugerahkan sukacita kepada saudara
berdua dan memberkati saudara dengan keturunan.}
\BU{Amin}
\BI{Semoga Putera Allah Yang Tunggal penuh kasih sayang membantu saudara dalam
untung dan malang}
\BU{Amin}
\BI{Semoga Roh Kudus selalu mencurahkan cinta kasihNya dalam hati saudara.}
\BU{Amin}
\BI{Dan semoga saudara sekalian diberkati oleh Allah yang mahakuasa, Bapa dan
Putera dan Roh Kudus}
\BU{Amin.}
\subjudul{Pengutusan}
\BI{Saudara sekalian, dengan ini perayaan Ekaristi dan penerimaan sakramen
perkawinan sudah selesai.}
\BU{Syukur kepada Allah}
\BI{Marilah kita memberi kesaksian tentang kebaikan Tuhan}
\BU{Amin.}

\subjudul{Persembahan kepada Bunda Maria}
\keterangan{Kedua mempelai meninggalkan altar menuju patung Bunda Maria membawa
bunga dan menghidupkan lilin}

\lagu{Lagu - \lagusungkemmaria}

\BPLW{Santa Maria, Bunda Yesus dan Bunda kami tercinta, hari ini merupakan hari
yang amat membahagiakan kami berdua. Di dalam iman akan Yesus Kristus, kami
berdua telah berjanji, untuk setia satu sama lain, dalam hidup bersama sebagai
suami istri. Kami sadar akan tanggung jawab yang harus kami emban dalam hidup
perkawinan kami, karena itu kami datang untuk memohon doa dan restumu ya Bunda
kami.

Doakanlah kami, agar cinta kami tahan uji dalam suka dan duka. Sertailah kami
dalam perjalanan keluarga kami. Dan semoga doamu mendatangkan berkat melimpah
bagi kami. Dan semoga kami dapat meneladan keluarga kudus Nasaret.

Salam Maria… 3$\times$}


\lagu{Lagu Penutup - \lagupenutup}

\newpage \pagestyle{empty}
\judul{UCAPAN TERIMA KASIH}
\noindent Dengan penuh syukur dan sukacita dalam kasih Tuhan, kami mengucapkan banyak
terima kasih kepada:
\large
\flushright{
\textbf{Romo \romo}\\
yang telah berkenan mempimpin perayaan ekaristi penerimaan sakramen perkawinan
ini.

\textbf{\saksisatu dan \saksidua}\\
yang telah berkenan menjadi saksi pernikahan ini.

\textbf{Paduan suara kelompok Fr. Fery}\\
yang telah menyemarakkan perayaan ekaristi ini.

\textbf{Segenap petugas liturgi}\\
(putra altar, lektor, petugas persembahan) \\
yang telah membantu kelancaran penerimaan sakramen pernikahan ini.

\textbf{Bapak Yoseph Samin, ketua Stasi Maguwo, dan Bapak Agung Dananjaya, ketua
lingkungan St. Petrus, serta segenap pengurus gereja}\\
yang berkenan mempersiapkan sarana pelaksanaan perayaan ekaristi penerimaan
sakramen pernikahan ini. 

\textbf{Segenap keluarga dan orang-orang terkasih}\\
yang telah berkenan hadir memberikan cinta, doa, dan restu dalam perayaan
ekaristi penerimaan sakramen pernikahan ini.

Semoga Tuhan memberkati dan memelihara ikatan kasih\\ di antara kita semua.

Amin.

\bigskip 

Kami yang berbahagia\\
Kel. \bpcamantri \\
Kel. \bpcamantra\\
Anto \& Zeli
}

\newpage \pagestyle{empty}
%\newsavebox\IBox
\sbox\IBox{\includegraphics{ornamen.eps}}
% set up the picture environment
\psset{unit=1in}
\begin{pspicture}(6.5in,6in)
% set up the fonts we use
\DeclareFixedFont{\PT}{T1}{ppl}{b}{it}{0.5in}
\DeclareFixedFont{\PTsmall}{T1}{ppl}{b}{it}{0.4in}
\DeclareFixedFont{\PTsmallest}{T1}{ppl}{b}{it}{0.3in}
\DeclareFixedFont{\PTtext}{T1}{ppl}{b}{it}{11pt}
\DeclareFixedFont{\Logo}{T1}{pbk}{m}{n}{0.3in}
% place the front cover picture
\rput[lm](0,2.5){\usebox\IBox}
\newsavebox\Blurbbox
\sbox\Blurbbox{\begin{minipage}{4.5in}
{\begin{flushright} \PTsmallest Ia membuat segala sesuatu indah pada waktunya,\\ 
bahkan Ia memberikan kekekalan dalam hati mereka.\\ {~} \\ {~} \\{~} \\ {~} \\
Indah, saat Ia\\ mempertemukan kami,\\
Indah, saat Ia\\ menumbuhkan cinta-kasih kami,\\
Indah pula, saat Ia \\menyatukan cinta-kasih kami.\\ {~} \\ {~} \\
Syukur kepada Tuhan atas kasih dan rencana-Nya yang indah dalam hidup kami\\

\end{flushright}}
\end{minipage}}
% And position the box
\rput[rm](6,2.5){\usebox\Blurbbox}% Then we close all open environments
\end{pspicture}


\end{document}
